\documentclass[11pt]{article}
\usepackage[margin=1in]{geometry}
\usepackage[T1]{fontenc}
\usepackage{amsmath,amssymb,booktabs,microtype}
\usepackage[hidelinks]{hyperref}
\setlength{\parindent}{0pt}
\setlength{\parskip}{0.6em}
\begin{document}
\begin{center}
{\Large\bfseries Reproducibility materials}\\[0.6em]
{\large Frequency Saturation and Certified Computation\\of Relaxation Functionals}\\[0.4em]
Haotian Zhong
\end{center}

\textbf{Archive:} \texttt{03\_Reproducibility\_Materials.zip}.
The archive contains the exact inputs, frozen certificates, checker source,
recorded seeds, and table-generation scripts supporting Section~7 of the
manuscript. All central proofs are contained in the manuscript. A file manifest
identifies the complete code and data snapshot by SHA-256 hashes.

\textbf{Quick start.} Extract the archive, open the resulting
\texttt{Supplementary\_Materials} directory, and execute
\begin{verbatim}
python3 run_all.py
\end{verbatim}
Python~3.10 or later suffices for frozen verification; the replay uses only the
standard library. A fresh run verifies the manifest and case inventory, runs the
checkers in disposable working directories, and regenerates the numerical tables.
The report is written to \texttt{replay\_output/SUMMARY.json}. Optional candidate
regeneration and its additional dependencies are described separately in
\texttt{REGENERATION.txt}.

\textbf{Evidence map.}
\begin{description}
\item[Sections 3--4; Section 7.1; Table 1; Figure 1.] Twenty-two certificates for
full-frequency positive pairs, four polynomial reproduction maps, and twelve
equal-effort design certificates. The frequency supremum is certified by an exact
response identity and a rational enclosure of its unique maximizer. Figure~1 is
drawn from the closed-form identity by the script
\texttt{make\_figure\_response\_difference.py} in the \texttt{figures} directory of
the manuscript source, after a cross-check against the archived enclosures.
\item[Proposition 4.6; Section 7.2.] Sixteen zero-noise identification
certificates, with positive spectra, identical original observations, and
nonzero target differences.
\item[Sections 5 and 6; Section 7.3 (Tables 2 and 3).] Eighteen stopped-cost
certificates and two executed stopping trajectories. The trajectories include six
stage inputs and twelve continuous endpoint certificates.
\item[Section 7.4 (Table 4).] Six strict-feasibility anchors on quadratic guard
grids, the relative-width checks, and table regeneration.
\end{description}

\textbf{Replay accounting.} The frozen suites comprise 80 certificate tasks and
154 named algebraic, mechanism, and invalid-input checks, including 53 intentional
rejections. The twelve continuous endpoints occur inside the two trajectory tasks.
These counts describe computational checks of the archived evidence. The two
seeded trajectories exercise the stopping algorithm; the coverage argument is
proved in Section~6.

\textbf{Provenance and navigation.} \texttt{SECTION\_MAP.txt} gives exact paths,
\texttt{INDEX.txt} lists the shipped files, and
\texttt{MANIFEST\_SHA256.json} records their hashes. Original suite names and
receipts retain their provenance. \texttt{TRANSFER\_NOTE.txt} explains the title
change and the metadata updates for the present submission. Copyright and
third-party dependency information appear in
\nolinkurl{RIGHTS_AND_DEPENDENCIES.txt}.
\end{document}
